\documentclass[11pt]{article}

% Standard packages
\usepackage[utf8]{inputenc}
\usepackage[T1]{fontenc}
\usepackage{amsmath,amssymb}
\usepackage{graphicx}
\usepackage{hyperref}
\usepackage{booktabs}
\usepackage{algorithm}
\usepackage{algpseudocode}
\usepackage[margin=1in]{geometry}
\usepackage{natbib}

% Title and authors
\title{Projections as Inductive Biases for Corpus-Based Language Models}

\author{
  Draft for Future Development
}

\date{\today}

\begin{document}

\maketitle

\begin{abstract}
\textbf{[PLACEHOLDER]} This paper explores how projection functions---transformations applied to queries or data---can serve as inductive biases in corpus-based language models. We present a unified framework where query projections ($\pi_Q$) and data projections ($\pi_D$) provide complementary mechanisms for improving generalization beyond exact string matching. We analyze the duality between runtime query transforms and corpus augmentation, discussing trade-offs in terms of storage, latency, and semantic coverage. This framework extends corpus-based models like Infinigram to handle case-insensitive matching, whitespace normalization, and other practical variations without modifying the core suffix array machinery.

\vspace{1em}
\noindent\textit{Note: This content was extracted from the Infinigram paper for separate development. See the main Infinigram paper for the core suffix array language modeling approach.}
\end{abstract}

\section{Introduction}

Corpus-based language models using suffix arrays \citep{manber1993suffix} provide exact pattern matching with $O(m \log n)$ query complexity. However, exact matching can be overly rigid: the query ``The Cat'' will not match corpus occurrences of ``the cat'' despite semantic equivalence for many applications.

This paper presents a \emph{projection framework} that addresses this limitation through controlled relaxation of exact matching. The key insight is that both query-time transformations and data augmentation can be viewed as \emph{projections} that map strings into equivalence classes. These projections function as \emph{inductive biases}---prior assumptions that improve generalization by treating syntactically different strings as semantically equivalent.

\subsection{Contributions}

Our contributions are:

\begin{enumerate}
    \item A unified \textbf{projection framework} with query projections ($\pi_Q$) and data projections ($\pi_D$) that formalize the relationship between runtime transforms and corpus augmentation.

    \item Analysis of \textbf{runtime query transforms} including case normalization, whitespace handling, and Unicode normalization.

    \item A formalization of \textbf{corpus augmentation as data projection}, showing how augmented corpora implement the same semantic relaxations with different trade-offs.

    \item A \textbf{duality theorem} establishing when query transforms and corpus augmentation produce equivalent results, and when they diverge.

    \item Discussion of \textbf{future directions} including semantic clustering, lemmatization, and embedding-based soft matching.
\end{enumerate}

\section{The Projection Framework}

We formalize projections as functions that map strings to equivalence classes, enabling matches between syntactically different but semantically equivalent strings.

\subsection{Query Projections}

A \textbf{query projection} $\pi_Q: \Sigma^* \to \Sigma^*$ transforms the query before matching against the corpus. Given a query $q$ and corpus $D$, we search for:
\begin{equation}
\text{matches}(\pi_Q(q), D) = \{p \in D : \pi_Q(q) \sqsubseteq D[p:]\}
\end{equation}
where $\sqsubseteq$ denotes prefix matching.

\textbf{Example.} With $\pi_Q = \texttt{lowercase}$, the query ``Hello World'' becomes ``hello world'', matching lowercase occurrences in the corpus.

\subsection{Data Projections}

A \textbf{data projection} $\pi_D: \Sigma^* \to \Sigma^*$ transforms the corpus at index time. Given query $q$ and projected corpus $\pi_D(D)$:
\begin{equation}
\text{matches}(q, \pi_D(D)) = \{p \in \pi_D(D) : q \sqsubseteq \pi_D(D)[p:]\}
\end{equation}

\textbf{Example.} With $\pi_D = \texttt{lowercase}$, the corpus ``Hello World'' is indexed as ``hello world'', matching lowercase queries.

\subsection{Symmetric Projections}

When $\pi_Q = \pi_D = \pi$, we have a \textbf{symmetric projection}:
\begin{equation}
\text{matches}(\pi(q), \pi(D))
\end{equation}

Symmetric projections define equivalence classes: strings $s_1$ and $s_2$ are equivalent iff $\pi(s_1) = \pi(s_2)$. This is the natural formulation for case-insensitive matching, where both query and corpus are lowercased.

\section{Runtime Query Transforms}

Infinigram supports several runtime query transforms that implement query projections without modifying the indexed corpus.

\subsection{Available Transforms}

\begin{description}
    \item[\texttt{lowercase}] Maps all characters to lowercase: ``Hello'' $\to$ ``hello''

    \item[\texttt{casefold}] Unicode-aware case folding for international text: ``Stra{\ss}e'' $\to$ ``strasse''

    \item[\texttt{strip}] Removes leading/trailing whitespace: ``\ hello\ '' $\to$ ``hello''

    \item[\texttt{normalize\_whitespace}] Collapses runs of whitespace to single spaces: ``hello\ \ \ world'' $\to$ ``hello world''
\end{description}

\subsection{Transform Composition}

Transforms can be composed: $\pi = \pi_2 \circ \pi_1$ applies $\pi_1$ first, then $\pi_2$. Common compositions include:

\begin{itemize}
    \item \texttt{strip $\circ$ lowercase}: Normalize case and remove edge whitespace
    \item \texttt{normalize\_whitespace $\circ$ lowercase}: Full text normalization
\end{itemize}

The order of composition matters in general, though many common transforms commute.

\subsection{Implementation}

Query transforms are applied in the \texttt{predict} method:

\begin{algorithmic}[1]
\Function{PredictWithTransform}{$q$, $\pi_Q$}
    \State $q' \gets \pi_Q(q)$
    \State \Return \Call{Predict}{$q'$}
\EndFunction
\end{algorithmic}

This adds negligible overhead since $\pi_Q$ operates only on the query (typically $<$100 bytes), not the corpus.

\section{Corpus Augmentation as Data Projection}

An alternative to query transforms is \textbf{corpus augmentation}: including transformed versions of documents in the index. This implements data projections.

\subsection{Augmentation Formula}

Given corpus $D = \{d_1, d_2, \ldots, d_n\}$ and augmentation functions $\{a_1, a_2, \ldots, a_k\}$, the augmented corpus is:
\begin{equation}
D' = D \cup \bigcup_{i=1}^{n} \bigcup_{j=1}^{k} \{a_j(d_i)\}
\end{equation}

Each document $d_i$ contributes $k+1$ versions: the original plus $k$ augmented variants.

\textbf{Example.} With $a_1 = \texttt{lowercase}$ and $a_2 = \texttt{uppercase}$:
\begin{align*}
D &= \{\text{``Hello World''}\} \\
D' &= \{\text{``Hello World''}, \text{``hello world''}, \text{``HELLO WORLD''}\}
\end{align*}

\subsection{Implementation}

Corpus augmentation is implemented in the \texttt{build\_corpus\_with\_augmentation} function:

\begin{algorithmic}[1]
\Function{BuildWithAugmentation}{$\text{docs}$, $\text{augmentations}$}
    \State $\text{corpus} \gets [\,]$
    \For{$d$ \textbf{in} $\text{docs}$}
        \State $\text{corpus}.\text{append}(d)$
        \For{$a$ \textbf{in} $\text{augmentations}$}
            \State $\text{corpus}.\text{append}(a(d))$
        \EndFor
    \EndFor
    \State \Return \Call{JoinWithSeparator}{$\text{corpus}$}
\EndFunction
\end{algorithmic}

\section{Duality and Trade-offs}

Query transforms and corpus augmentation provide dual mechanisms for achieving similar goals. This section analyzes their relationship and trade-offs.

\subsection{Equivalence Conditions}

For a symmetric projection $\pi$, query transform and corpus augmentation produce equivalent results when:
\begin{enumerate}
    \item The corpus is \emph{exhaustively augmented}: $D' = \bigcup_{s \in \pi^{-1}(\pi(D))} \{s\}$
    \item The query matches \emph{all} augmented variants equally
\end{enumerate}

In practice, exhaustive augmentation is infeasible, so the two approaches diverge.

\subsection{Trade-off Comparison}

\begin{table}[h]
\centering
\begin{tabular}{lcc}
\toprule
\textbf{Aspect} & \textbf{Query Transform} & \textbf{Corpus Augmentation} \\
\midrule
Storage & $1\times$ (no increase) & $(k+1)\times$ per augmentation \\
Index time & None & $O(k \cdot n \log n)$ \\
Query latency & $O(1)$ transform overhead & None (pre-computed) \\
Coverage & Query-side only & Full bidirectional \\
Flexibility & Dynamic (per-query) & Static (build-time) \\
\bottomrule
\end{tabular}
\caption{Trade-offs between query transforms and corpus augmentation.}
\label{tab:tradeoffs}
\end{table}

\textbf{Query transforms} are preferred when:
\begin{itemize}
    \item Storage is constrained
    \item Different queries need different transforms
    \item The corpus is frequently updated (re-indexing is expensive)
\end{itemize}

\textbf{Corpus augmentation} is preferred when:
\begin{itemize}
    \item Query latency is critical
    \item Bidirectional matching is required (corpus and query both vary)
    \item The augmentation set is small and fixed
\end{itemize}

\subsection{Hybrid Approaches}

Practical systems often combine both approaches:
\begin{itemize}
    \item Apply cheap transforms (whitespace normalization) at query time
    \item Pre-compute expensive transforms (lemmatization) in the corpus
    \item Use augmentation for high-frequency variations (common misspellings)
\end{itemize}

\section{Theoretical Perspective}

The projection framework connects corpus-based models to neural language models through the concept of \emph{inductive bias}.

\subsection{Projections as Inductive Biases}

An \textbf{inductive bias} is a set of assumptions that a learning algorithm uses to generalize beyond observed data. In neural networks, architectural choices (convolutions, attention, recurrence) implement implicit inductive biases about the structure of the target function.

Projections in corpus-based models serve an analogous role:
\begin{itemize}
    \item \textbf{Case insensitivity} assumes semantic equivalence across case variations
    \item \textbf{Whitespace normalization} assumes formatting variations are noise
    \item \textbf{Lemmatization} assumes morphological variants share meaning
\end{itemize}

Unlike neural inductive biases, projection-based biases are:
\begin{enumerate}
    \item \textbf{Explicit}: The bias is a known, inspectable function
    \item \textbf{Composable}: Multiple biases combine predictably
    \item \textbf{Controllable}: Biases can be enabled/disabled per-query
\end{enumerate}

\subsection{Connection to Smoothing}

Traditional n-gram smoothing (Kneser-Ney, Witten-Bell) can be viewed as implementing a projection that maps unseen n-grams to ``similar'' seen n-grams. The smoothing distribution defines an implicit similarity metric.

Our projection framework makes this explicit: rather than statistical interpolation, we specify exactly which strings should be treated as equivalent.

\section{Future Directions}

The projection framework opens several research directions:

\subsection{Semantic Clustering}

Map tokens to \emph{concept IDs} based on distributional similarity:
\begin{equation}
\pi_{\text{cluster}}(\text{``cat''}) = \pi_{\text{cluster}}(\text{``kitten''}) = \text{FELINE}
\end{equation}

This requires pre-computed embeddings but enables soft semantic matching within the exact-match framework.

\subsection{Lemmatization}

Map words to root forms:
\begin{equation}
\pi_{\text{lemma}}(\text{``running''}) = \pi_{\text{lemma}}(\text{``ran''}) = \text{``run''}
\end{equation}

This handles morphological variation but requires language-specific stemmers or lemmatizers.

\subsection{Embedding-Based Soft Matching}

Use neural embeddings to define continuous projections:
\begin{equation}
\pi_{\text{embed}}(s) = \text{argmin}_{c \in \mathcal{C}} \|E(s) - E(c)\|
\end{equation}

where $\mathcal{C}$ is a codebook and $E$ is an embedding function. This bridges the gap between exact corpus matching and neural semantic similarity.

\subsection{Learned Projections}

Train projections end-to-end to optimize downstream task performance:
\begin{equation}
\pi^* = \text{argmin}_\pi \mathcal{L}(\text{LM}_{\pi(D)}, \text{eval data})
\end{equation}

This could automatically discover useful equivalence classes from data.

\section{Conclusion}

The projection framework provides a principled approach to relaxing exact matching in corpus-based language models. By viewing query transforms and corpus augmentation as dual mechanisms implementing the same underlying concept---projection to equivalence classes---we gain clarity about their trade-offs and can design hybrid systems that combine their strengths.

The connection to inductive biases in neural networks suggests that corpus-based models, despite their simplicity, can incorporate sophisticated prior knowledge through carefully designed projections. Future work on semantic clustering, lemmatization, and embedding-based matching will further extend the expressive power of this approach.

\bibliographystyle{plainnat}
\bibliography{references}

\end{document}
